\begin{figure*}[t]
\centering
\begin{tikzpicture}[
  font=\small,
  box/.style={rounded corners=2pt,draw=black!30,thick},
  lbl/.style={font=\footnotesize\bfseries,text=cInk},
]
% ===== 左：光滑物体+缺陷 =====
\node[lbl] at (1.6,3.3) {(a) 光滑物体 + 缺陷};
% input
\node[box,fill=cBase!35,minimum width=2.4cm,minimum height=1.5cm] (ia) at (1.6,2.2) {};
\fill[cBad] (1.35,2.35) rectangle (2.05,2.5);
\node[font=\scriptsize,text=black!55] at (1.6,1.25) {输入};
% HF response
\node[box,fill=cFreq!22,minimum width=2.4cm,minimum height=1.5cm] (ha) at (1.6,-0.6) {};
\fill[cFreq] (1.35,-0.45) rectangle (2.05,-0.3);
\node[font=\scriptsize,text=cFreq!60!black] at (1.6,-1.55) {缺陷凸显（好）};

% ===== 右：粗糙正常纹理 =====
\node[lbl] at (6.0,3.3) {(b) 粗糙材质，正常};
\node[box,fill=cBase!35,minimum width=2.4cm,minimum height=1.5cm] (ib) at (6.0,2.2) {};
\foreach \y in {1.65,1.85,2.05,2.25,2.45,2.65}{\draw[black!35,line width=0.5pt] (4.85,\y)--(7.15,\y);}
\node[font=\scriptsize,text=black!55] at (6.0,1.25) {无缺陷};
\node[box,fill=cFreq!22,minimum width=2.4cm,minimum height=1.5cm] (hb) at (6.0,-0.6) {};
\foreach \y in {-1.15,-0.95,-0.75,-0.55,-0.35,-0.15}{\draw[cFreq!80!black,line width=0.6pt,opacity=0.75] (4.85,\y)--(7.15,\y);}
\node[font=\scriptsize,text=cGlobal!45!black] at (6.0,-1.55) {整片变亮（坏）};

% 行标题
\node[lbl,rotate=90,text=cFreq!60!black] at (-0.35,-0.6) {CLIP 高频响应};
\node[lbl,rotate=90,text=black!55] at (-0.35,2.2) {输入 / 真值};

% 箭头
\draw[-{Latex[length=2mm]},black!55] (ia) -- (ha);
\draw[-{Latex[length=2mm]},black!55] (ib) -- (hb);

% ===== 中间洞察框 =====
\node[box,fill=cGlobal!18,draw=cGlobal!70,minimum width=3.2cm,align=center,font=\footnotesize] (ins) at (9.9,1.7)
{\textbf{相同的高频幅度}\\[1pt] 在 (a) 是异常\\ 在 (b) 是正常\\[2pt]\textbf{$\Rightarrow$ 固定阈值失效}};

% ===== 解法框 =====
\node[box,fill=cOurs!18,draw=cOurs!70,minimum width=6.6cm,align=center,font=\footnotesize] (sol) at (11.6,-0.9)
{\textbf{本文思路：} 估计每张图\emph{自己}的正常高频水平，只标记超出该逐图参照的部分};

\draw[-{Latex[length=2mm]},cOurs!70,thick] (ins) -- (sol);
\end{tikzpicture}
\caption{CLIP patch 特征的高频响应对缺陷（a）敏感，但对正常粗糙材质（b）同样敏感。单一固定阈值无法区分“异常”与“材质”。我们的方法估计每张图自己的正常高频水平，只标记超出该逐图参照的部分。}
\label{fig:motivation}
\end{figure*}
